\section{Introduction and design position}
\label{sec:position}

\subsection{What \edet\ is}

\edet\ is a clearing medium for a community of mutual producers: people and small
enterprises who buy from one another regularly enough to clear their obligations against
each other instead of settling in national currency.

Its unit is a debt contract, created by a sale from buyer to seller in a community unit.
It discharges when the buyer delivers value onward --- to the seller, or to anyone the
seller directs it to --- so an obligation moves through the community as production does
and is worked off rather than paid out of a stock. There is no token, no bearer
instrument, no balance and nothing to hoard: a large claim is a large expectation of
future delivery, not a holding. What a member accumulates is \emph{capacity}, a ceiling on
how much the community will underwrite them for: bounded, non-transferable, not for sale,
living in the stakes others have placed on them and falling when they withdraw
(\S\ref{sec:standing}).

The design is measured for communities up to the order of $10^5$ accounts on one ledger
(\S\ref{sec:implementation}); a larger ledger is an open problem this paper states rather
than escapes, since a stake graph cannot be split without severing every edge across the
cut.

\subsection{The two structural commitments}

\paragraph{Capacity is a minimum cut.} A member's capacity is the maximum flow from the
community's underwriters to their account, across directed stakes recording what creditors
have placed. Nobody can stake beyond what they may confer, outstanding credit reserves the
flow that justified it, and an underwriter is any member who has declared a supply they
will stand behind. The ledger stores stakes and outstanding obligations; capacity is the
cut computed from them.

One equation settles four questions a credit system normally answers with policy.
Admission: an account with no stakes has capacity zero, so creating one grants nothing.
Reputation: the stakes \emph{are} the reputation. Sybil resistance: the capacity of any set
of accounts is bounded by the stakes crossing into it, so identities are free and worthless
together. Fraud: volume internal to a set never crosses that set's own boundary, so
fabricated trade is arithmetically void. These are theorems rather than process ---
Theorem~\ref{thm:cut} is argued once --- where a rule that must be remembered can be
forgotten, one that must be enforced can be evaded, and a bound re-argued for each new
attack is eventually re-argued wrongly.

\paragraph{Capacity underwrites, it does not permit.} Credit within a debtor's capacity is
\emph{insured}: it reserves flow, and the community's recourse machinery stands behind it.
Credit beyond it is permitted but \emph{uninsured}, and the creditor bears it alone. This
is what lets a community begin: first trades are uninsured, they settle, they create
stakes, and capacity is their residue (\S\ref{sec:standing}).

\paragraph{Ledger-scoped total order.} A ledger is one replicated state machine under
Byzantine-fault-tolerant consensus among validators its members operate themselves, with
no world-global chain, no gas and no fee token, and every quantity in the model computed
within one order. A ledger is not a community: a community is a region of the stake
graph, one ledger carries as many as its members form, and communities that expect to
trade belong on one, which is why there is no bridge (\S\ref{sec:model},
\S\ref{sec:oneledger}). Consensus is load-bearing: reservation makes capacity a conserved
budget over a shared graph, and a conserved budget needs a serialization point
(Remark~\ref{rem:order}).

\subsection{The boundaries that remain}

Two limits are structural, and neither is acceptance. The limitation usually written here
--- that nobody can be made to want a medium they cannot hold --- describes a bearer
instrument; a sale here discharges before it creates (\S\ref{sec:sale}), so a producer
with obligations is paid by ceasing to owe. What binds is the thesis's antecedent, that
enough of what a community buys is available from within it: a property of the community,
measurable before founding, which no protocol changes (\S\ref{sec:adoption}). And zero is
absorbing: a community that underwrites nothing at genesis can never leave that state from
inside, and \S\ref{sec:seed-amendments} is its only door. The newcomer's first uninsured
trade is not a third limit but a priced cost: somebody extends a small uninsured risk, it
settles, and the stake it writes makes the next trade insured (Remark~\ref{rem:friction}).

\subsection{Reader's map}

\S\ref{sec:model} fixes the system model; \S\ref{sec:medium} defines the contract and
how it clears; \S\ref{sec:standing} defines capacity and is the core of the design;
\S\ref{sec:stability} treats containment; \S\ref{sec:recourse} covers default,
substitution and arbitration; \S\ref{sec:governance} covers governed constants and
re-denomination; \S\ref{sec:security} states and proves the security results;
\S\ref{sec:economics} through \S\ref{sec:provision} take the economic position; and
\S\ref{sec:implementation} and \S\ref{sec:verification} describe what is built and what
is checked.
